\section{Discussion}\label{sec:discussion}

We have shown that the compressive sector of KRMHD is a continuum of passive scalars, the slices, labelled by the parallel phase velocity. Without collisions, each slice is carried by the $\boldsymbol E\times\boldsymbol B$ flow and streams along the perturbed magnetic field at its own speed, independently of all other slices, in three dimensions and for any smooth Alfv\'enic fields (theorem~\ref{thm:slice}). The Landau coupling, which links the distribution at different velocities, is absorbed into the definition of the slices because its strength does not depend on the wavenumber, and the slices are passive because the Alfv\'enic fields do not depend on the compressive distribution. A Hermite hierarchy truncated at order $M$ has the same structure, with $M+1$ slices at the eigenvalues of its ladder matrix (theorem~\ref{thm:trunc}), which for the closure $g_{M+1}=0$ are the nodes of the Gauss rule of the Case--Van Kampen spectral measure (\ref{eq:gauss}). With the decomposition we have derived a family of invariants that contains the free energy $W$ and the invariants $\Gamma$ and $\Hph$ of CMM, described phase mixing as the decorrelation of slices, and identified the two slices, at $v_\parallel=\pm v_A$, that are advected by a single Elsasser field. We have also proved the frozen-line theorem, which builds on the Lagrangian-frame linearity of S09 and \citet{Kunz2015} and gives exact statements about volume-integrated statistics when the induction equation is ideal. A symmetry that exchanges the slices at $u$ and $-u$ gives a theorem on forced ensembles that are mirror images of each other.

These results hold in the limit $k_\perp\rho_i\to0$. They apply to each compressive branch separately and to the Laguerre sectors of order $l\ge2$, all of which have $\Lambda\notin[0,1]$ (section~\ref{sec:model}). They concern the collisionless equations without forcing or spatial dissipation of the compressive distribution, except where collisions, dissipation or forcing are stated explicitly. The slice theorem, the invariants, the two-slice law and the representations of the Hermite flux and spectrum need only smooth Alfv\'enic fields, which may be forced, resistive or prescribed. The frozen-line theorem and its corollaries need in addition the hypotheses (\ref{eq:hyp}): an ideal induction equation, no spatial dissipation of the compressive distribution, and Alfv\'enic fields that remain smooth on the interval considered. They allow collisions that act at each point in space separately, and forcing of the compressive distribution that is white in time, for which they concern the means of quadratic statistics.

\subsection{Velocity-space transfer and echoes}\label{dis:transfer}

Without collisions, free energy does not move between phase velocities, because the energy of every slice is conserved. It moves between Hermite moments only when the correlations between different slices change, and these change only through the relative streaming of the slices along the perturbed field, at a rate proportional to the difference of their speeds (proposition~\ref{prop:twoslice}). Phase mixing is the decorrelation of slices and phase unmixing is their recorrelation. Neither definition refers to a wavenumber, and both apply when the field lines move. Velocity-space transfer in collisionless KRMHD is therefore the decorrelation and recorrelation of independent passive scalars.

The invariants found here for the compressive sector of collisionless KRMHD are the family $\Qphi$, which contains $W$, $\Gamma$ and $\Hph$, the Casimirs of each slice and the invariants that couple different sectors. All of them are functionals of the statistics of single slices, or of slices of different sectors at the same $u$. They fix the distribution of free energy over phase velocity but not the correlations between different slices of one sector. Linear phase mixing in a uniform field is an exact solution of KRMHD that conserves all of them while it moves all of the free energy, and all of $\Gamma$, to arbitrarily large $m$, so the invariants impose no constraint of the kind that the argument of \citet{Fjortoft1953} gives in two-dimensional hydrodynamics (proposition~\ref{prop:nofjortoft}). This completes the heuristic argument of CMM for $\Gamma$ and extends it to every invariant of collisionless KRMHD that depends on the compressive distribution alone. The argument concerns exact transient solutions. A version for the forced and dissipative steady states, to which the cascade criterion of CMM refers, is not attempted here.

Under the hypotheses (\ref{eq:hyp}) the turbulent advection does not convert phase mixing into unmixing in the Hermite spectrum. The spectrum $W(m,t)$ and the field-aligned Hermite flux evolve as in linear phase mixing along the field lines frozen at the initial time (corollary~\ref{cor:singlepoint}), and with forcing that is white in time their means are sums of such linear responses, each along the field lines frozen at its injection time (corollary~\ref{cor:forcing}). Stochastic echoes in $W(m)$ therefore require a departure from these hypotheses: a non-ideal term in the induction equation, spatial dissipation of the compressive distribution, forcing of the distribution that is correlated in time, or, in the limit of vanishing resistivity, a failure of deterministic flux freezing for fields that do not remain smooth \citep{Eyink2013,Eyink2018,Jafari2026}. The KRMHD simulations in which echoes were reported \citep{Meyrand2019} dissipate both the magnetic field and the compressive distribution in space, so the theorem does not contradict them. It shows that an echo of the kind they describe, in which the sign of the wavenumber along the field is reversed, requires these terms, or a failure of flux freezing in the limit in which they vanish. Their Hermite flux resolved in $k_\perp$ shells is not a single-point statistic, and the theorem constrains only its sum over shells (section~\ref{sec:lagrangian}). The gyrokinetic simulations of \citet{Kawazura2020} retain finite-$k_\perp\rho_i$ effects even in their RMHD range, so the theorem does not apply to them directly.

\subsection{Simulations and diagnostics}\label{dis:numerics}

Several practical points follow for simulations and their diagnostics. Diagnostics of velocity-space transfer should use the derivative along the perturbed field. The Hermite flux computed with $\partial_z$ alone omits the field-line term $\{\Psi,\cdot\}$, which in the strongly nonlinear random states of our checks was between 29\% and 253\% of the field-aligned flux in magnitude, depending on $m$ (section~\ref{sec:decorrelation}). A separation of phase mixing from unmixing by the sign of $k_z$ misreads phase mixing along tilted field lines. For streaming along uniformly tilted field lines and a Gaussian ensemble of modes with effective tilt angle $\alpha$, the net flux computed with $k_z$ is $\cos^2\alpha$ times the field-aligned flux, and a fraction $\alpha/\pi$ of the spectral weight lies in modes whose $k_z$ and field-aligned wavenumber have opposite signs, on which the split exchanges phase mixing and unmixing (\ref{eq:tiltcos}). In critically balanced turbulence $\tan\alpha$ is of order unity, so this estimate suggests a misattribution of order unity there; it has been derived for uniformly tilted field lines, not measured in turbulence. In slice variables the field-aligned flux is the Christoffel--Darboux transform of the rate of change of the slice correlations (\ref{eq:fluxCD}), a form that refers to no wavenumber.

Diagnostics of the phase-space helicity should use an odd truncation order $M$, for which, with the closure $g_{M+1}=0$, the conserved discrete helicity is CMM's $\Hph$ evaluated on the retained moments. At even $M$, $u=0$ is an eigenvalue of the truncated ladder matrix, the restriction of $\Hph$ to the retained moments is not conserved, and the conserved principal-value form, which omits the slice at $u=0$, converges to $\Hph$ only as $M^{-1/2}$, with errors of 5--18\% at $M=64$ and $128$ for a shifted Maxwellian (\ref{eq:Hoddeven}).

In a Hermite code that treats the nonlinear terms explicitly, slice $j$ is advected with the stream function $\Phi+su_j\Psi$, so the time step is limited by the fastest slice, whose speed involves the largest eigenvalue $u_{\max}$ of the truncated ladder matrix; $u_{\max}=10.6$ at $M=64$ and $15.4$ at $M=128$ for the closure $g_{M+1}=0$, and $11.0$ and $15.7$ for the copy closure $g_{M+1}=g_{M-1}$ (\ref{eq:umax}). Estimates of the time step based on $\max(v_A,|\nabla_\perp z^\pm|,s|\nabla_\perp\Psi|)$, such as GANDALF's, omit this factor. The GX code uses the largest Gauss--Hermite node in its time-step estimate for parallel streaming \citep{Mandell2024}, but the speed $su_{\max}|\nabla_\perp\Psi|$ of the field-line part of the nonlinearity does not enter that estimate. For pure advection, a fourth-order Runge--Kutta treatment of the nonlinear terms is stable when $\Delta t\,\rho\le2\sqrt2$, where $\rho$ is the spectral radius of the advection operator of the fastest slice, while the explicit stages of the IMEX-RK222 scheme have no stable interval on the imaginary axis, so that their stability rests on dissipation (\ref{eq:explicitstability}).

Both closures of a truncated hierarchy are conservative, since no free energy leaves the hierarchy at $m=M$, so a phase-mixing packet that reaches $m=M$ is reflected. With hypercollisions its returning amplitude is reduced by the factor $e^{-X}$ of (\ref{eq:reflection}), and in a steady state forced by white noise in $g_0$ the reflected energy raises the Hermite spectrum at $m\ll M$ by the factor $\coth X$ above the spectrum without reflection (section~\ref{sec:truncated}). Both results were derived and tested for linear problems at a single wavenumber.

The decomposition also suggests a way to integrate the compressive sector. Each slice could be advanced as an independent passive scalar, for example by a semi-Lagrangian method along the characteristics of its velocity $\boldsymbol U_u$ (\ref{eq:Uu}), which could remove the stability restriction set by $u_{\max}$. The slices could be advanced in parallel, and only a collision operator that depends on velocity, through the kernel in (\ref{eq:sliceforcing}), would couple them. We have not implemented or tested such a scheme.

\subsection{Open questions}\label{dis:open}

Hypothesis (H2) of (\ref{eq:hyp}) excludes spatial dissipation of the compressive distribution. Each slice is a passive scalar, and in turbulent flows the rate at which the variance of a passive scalar is dissipated remains finite as the diffusivity goes to zero, which is the dissipative anomaly of passive-scalar mixing \citep{ShraimanSiggia2000,FalkovichGawedzkiVergassola2001}. The $k_\perp$ sink of each slice therefore need not disappear as the dissipation coefficients go to zero, and in turbulence its competition with phase mixing, rather than an ideal echo, may set the Hermite spectrum at large $m$; whether it does is an open question. A related question is whether the departures from the frozen-line predictions that spatial dissipation of the compressive distribution produces in our tests (sections~\ref{sec:lagrangian} and~\ref{sec:symmetry}), among them the departure from the parity lemma (lemma~\ref{lem:parity}), survive as the dissipation coefficients tend to zero; at fixed coefficients they do not decrease with resolution.

At finite $k_\perp\rho_i$ the Landau coupling depends on the wavenumber, the advection depends on $v_\perp$, and the compressive fluctuations act back on the Alfv\'enic ones (section~\ref{sec:invariants}), so the slices no longer decouple. For advection by a prescribed potential we have characterised the quadratic forms that remain conserved, but the field equations can make further forms invariant, as they do for the gyrokinetic free energy. CMM report the free energy and the gyrokinetic helicity as invariants of full gyrokinetics, and $\Gamma$ only for KRMHD. Whether $\Gamma$ has a corrected form that is conserved at order $k_\perp\rho_i$ has not been computed. In an electrostatic gyrokinetic model with Boltzmann electrons its analogue is broken at order $(k_\perp\rho_i)^2$, but that model lacks the electromagnetic terms of order $k_\perp\rho_i$.

Collision operators that depend on velocity are the only terms that exchange free energy between slices. For operators that are diagonal in the Hermite basis, such as the Lenard--Bernstein operator in $v_\parallel$ and the hypercollisions of Hermite codes, the coupling kernel is explicit (\ref{eq:sliceforcing}). How this exchange combines with the decorrelation of slices in turbulence is left open.

The frozen-line theorem holds on intervals on which the Alfv\'enic fields stay smooth. \citet{Eyink2013} reported a breakdown of flux freezing in high-conductivity magnetohydrodynamic turbulence, \citet{Eyink2018} argued that flux freezing fails in a turbulent inertial range, and \citet{Jafari2026} argued that for the rough fields of such a range the field lines remain stochastic in the limit of vanishing resistivity, so that flux is conserved only statistically (spontaneous stochasticity). Whether this happens in KRMHD turbulence, and how a stochastic relabelling of the slices would change the Hermite spectrum, are open questions.

Finally, the forcing results of section~\ref{sec:lagrangian} concern means over white forcing on finite intervals, and proposition~\ref{prop:nofjortoft} concerns transient solutions. Statistically steady states of forced and dissipative KRMHD turbulence, in which questions about cascades in velocity space are usually posed, were not treated here.

% Acknowledgements and any AI-use disclosure to be written by the author.
